\documentclass[10pt]{article}
\usepackage[letterpaper,margin=.85in]{geometry}
\usepackage[T1]{fontenc}
\usepackage{lmodern,microtype,amsmath,amssymb,amsthm,booktabs}
\usepackage[hidelinks,pdftitle={Cayley equivalence for short patterns},pdfauthor={1200}]{hyperref}
\newtheorem{theorem}{Theorem}
\setlength{\parindent}{0pt}
\setlength{\parskip}{.48em}
\raggedbottom
\title{\textbf{Cayley equivalence for short patterns}}
\author{1200}
\date{November 2026}
\begin{document}
\maketitle
\begin{abstract}
We verify Conjecture 8.6 of Cerbai, Claesson, Ernst and Golab for all packed
patterns of length at most five and for patterns of length six on at most
two letters. In these ranges, total Cayley avoidance counts through length
eleven already determine all fixed-alphabet avoidance counts. The proof
combines exact finite distinguishing counts with the published word-Wilf
equivalences of Jel\'{\i}nek and Mansour.
\end{abstract}

\section{Statement and reduction}
A packed word has a set of letters of the form $[k]=\{1,\ldots,k\}$;
the empty word has $k=0$. An occurrence of a packed pattern is a subsequence
with the same order relations and equalities. For a pattern $p$, let
$A_p(n,k)$ count the avoiding packed words of length $n$ on exactly $[k]$,
let $A_p(n)=\sum_k A_p(n,k)$, and let $W_p(n,K)$ count avoiding words on
$[K]$ without requiring every letter to occur.

\begin{theorem}
Suppose $p$ and $q$ have the same positive length $m$. If $m\le5$, or if
$m=6$ and both patterns use at most two distinct values, then the
following conditions are equivalent:
\begin{enumerate}
\setlength{\itemsep}{.14em}
\item $A_p(n)=A_q(n)$ for every $0\le n\le11$;
\item $A_p(n)=A_q(n)$ for every $n\ge0$;
\item $A_p(n,k)=A_q(n,k)$ for every $n,k\ge0$;
\item $W_p(n,K)=W_q(n,K)$ for every $n,K\ge0$.
\end{enumerate}
\end{theorem}

Split a word on $[K]$ according to its set of used letters. Increasing
relabelling preserves containment, so for each $n$,
\[
 W_p(n,K)=\sum_{k=0}^K \binom Kk A_p(n,k).
\]
This identity and binomial inversion give (3)$\Leftrightarrow$(4);
summing gives (3)$\Rightarrow$(2), and (2)$\Rightarrow$(1) is immediate.
We prove (1)$\Rightarrow$(4) by showing that distinct word-Wilf classes
have distinct finite total signatures.

\section{Class distinctions}
Use the word-Wilf classification in \cite[\S\S1.1, 1.4 and Appendix A]{jm}.
The authors' table \cite{table} lists representatives up to reversal and
rank complement; both operations preserve fixed-alphabet avoidance.
Expanding its rows covers exactly the 75 patterns of length four
(16 classes), 541 of length five (116 classes), and 63 at-most-binary
patterns of length six (18 classes), without overlap. Membership in
the same row, after expansion, supplies the all-length implication.

There is one pattern of length one. At length two the classes represented
by $11$ and $12$ have totals 6 and 4 at host length three. At length three
the classes represented by $111,112,123$ have totals 66, 60, and 56
at host length four. At pattern lengths four and five, direct enumeration
at host lengths eight and nine, respectively, separates all but the
corresponding rows listed below. At pattern length six with at most two
values, the vectors of totals at host lengths six through ten separate
all but the last two rows. Each row of this table supplies a further
host length that separates its pair.
\begin{center}
\begin{tabular}{clcrr}
\toprule
$m$ & Representatives $(p,q)$ & $n$ & $A_p(n)$ & $A_q(n)$\\
\midrule
4 & $1223,\ 1332$ & 9 & 4,472,716 & 4,472,712\\
5 & $11212,\ 12121$ & 10 & 98,890,988 & 98,890,980\\
5 & $11223,\ 11332$ & 11 & 1,445,533,180 & 1,445,533,200\\
5 & $12123,\ 13232$ & 10 & 93,909,363 & 93,909,357\\
5 & $12234,\ 12443$ & 10 & 87,548,792 & 87,548,796\\
5 & $13442,\ 14233$ & 10 & 87,410,614 & 87,410,626\\
6 & $112112,\ 121112$ & 11 & 1,616,420,895 & 1,616,420,896\\
6 & $111212,\ 112121$ & 11 & 1,616,424,886 & 1,616,424,885\\
\bottomrule
\end{tabular}
\end{center}
Thus (1) forces membership in the same word-Wilf class, which gives (4)
by the cited classification. This proves the stated part of Conjecture
8.6 \cite[p.~295]{cceg}. The original conjecture for arbitrary pattern
length is outside the result.

\section{Exact counts}
For the two complete censuses, enumerate each host on its actual
alphabet, retaining only words that use every value of that alphabet.
At a leaf, standardize each increasing selection of positions and count
each resulting pattern once for that host. The code checks the total
population against the sum of the layers given by the surjection recurrence
\[
 F(n,k)=k\bigl(F(n-1,k)+F(n-1,k-1)\bigr), \qquad F(0,0)=1.
\]
Deleting the last letter either leaves all $k$ letters or removes the
only instance of that letter, giving the two terms. Subtracting containing
hosts from the population gives avoidance.

For individual patterns, an alternative count starts at the pattern
and inserts one letter at every step. Either the letter is already
present, or it is a new rank and all old ranks at least as large are
increased by one. In any larger containing word, deleting a position
outside a fixed occurrence and standardizing gives a containing parent.
Conversely either insertion preserves that occurrence. This proves
coverage, and storing the words in sets handles repeated generation.

The large comparison $11223/11332$ uses a streaming recognizer.
Both patterns have the form $(a,a,d,d,x)$, with $a<d<x$ for the first
or $a<x<d$ for the second. While reading a prefix, maintain the minimum
repeated value, the least $a<d$ for each realized prefix $(a,a,d)$,
and the possible final ranks following any $(a,a,d,d)$. A new letter is
tested as a final rank before updating shorter prefixes. The least $a$
for each $d$ gives the whole interval of possible middle ranks; for the
other pattern, only the minimum completed $d$ matters. If both patterns
have occurred, their common remaining packed completions can be added
by a recurrence on the remaining positions and unused letters.
The streaming counts agree with insertion enumeration at lengths
nine and ten. The accompanying archive includes source code, execution
records, and a coverage checker.

\textbf{AI contributions: GLM5.3F and GPT Astra.}
GLM5.3F contributed to the source research. GPT Astra assisted with the
short-pattern classification, proof development, and independent verification.

\begin{thebibliography}{9}
\bibitem{cceg} G. Cerbai, A. Claesson, D. C. Ernst and H. Golab,
Pattern-avoiding Cayley permutations via combinatorial species,
\href{https://ajc.maths.uq.edu.au/pdf/94/ajc_v94_p267.pdf}
{Australas. J. Combin. \textbf{94}(2) (2026), 267--299}.
\bibitem{jm} V. Jel\'{\i}nek and T. Mansour, Wilf-equivalence on $k$-ary words,
compositions, and parking functions,
\href{https://doi.org/10.37236/147}{Electron. J. Combin. \textbf{16}(1) (2009), \#R58}.
\bibitem{table} V. Jel\'{\i}nek and T. Mansour,
\href{https://math.haifa.ac.il/toufik/enumerative/tablewords456.html}
{Supplementary classification table for patterns of lengths four, five, and six}.
\end{thebibliography}
\end{document}
